problem:  
Let \(M^k \subset S^{n}(1)\) be a compact **minimal submanifold** of the unit sphere with nontrivial normal bundle.  
For each \(v \in S^{n}(1)\), consider the **height function** \(f_v(x) = \langle x, v \rangle\) on \(M\).  For generic \(v\), \(f_v\) is a Morse function.  Denote by  
\[
\mu_i(M,v) = \#\{p \in M : \nabla f_v(p)=0,\ \operatorname{ind}_p(f_v)=i\}
\]
the number of critical points with Morse index \(i\).  

Assume \(M\) admits a smooth one-parameter family of normal deformations \(M_\varepsilon\), where \(M_0=M\), such that \(M_\varepsilon\) is minimal up to order \(O(\varepsilon^2)\) with respect to the induced metric from \(S^n(1)\).  Define the **Morse index distribution functional**
\[
\mathcal{F}_i(\varepsilon) = \int_{S^n} \mu_i(M_\varepsilon,v)\, d\sigma(v),
\]
the expected number of critical points of index \(i\) for random height directions.

**Conjectural problem:**  
Can the functions \(\mathcal{F}_i(\varepsilon)\) vary nontrivially with \(\varepsilon\) even when each \(M_\varepsilon\) remains minimal to second order?  
Equivalently, does there exist a compact minimal submanifold \(M\subset S^n(1)\) that admits an infinitesimal normal deformation preserving minimality up to \(O(\varepsilon^2)\) but producing a strictly first-order change in the averaged critical point statistics \(\{\mathcal{F}_i(\varepsilon)\}\)?  

In other words, can infinitesimal **symmetry breaking** in near-minimal deformations leave classical geometric invariants unchanged while still altering the *probabilistic Morse-theoretic profile* of height functions?

Why is it a "good" problem:  
This revised problem avoids the topological invariance obstruction that was present in the Morse–Euler sum and instead focuses on the finer distributional structure of Morse indices, which is not topologically fixed.  The question probes whether averaged Morse index statistics are **sensitive to infinitesimal geometric symmetry breaking** in minimal settings.  
It lies at the nexus of **Morse theory**, **integral geometry**, and **variational analysis of minimal submanifolds**.  A positive answer would uncover new quantitative invariants bridging stochastic topology and extrinsic curvature geometry; a negative one would reveal rigidity phenomena for Morse index distributions under near-minimal perturbations.  
The problem is concrete yet deep—its resolution requires unifying local curvature analysis with global average topological data—and it could potentially introduce a novel geometric statistic for detecting higher-order symmetry loss, making it an authentically “good” and forward-looking research-level problem.